% !TEX program = lualatex
% Yuho — research paper. arXiv preprint, attributed.
% Build with: make paper

\documentclass[sigconf,nonacm,screen]{acmart}

% --------------------------------------------------------------------
% Author macros (placeholders + \todo for skeleton sections)
% --------------------------------------------------------------------
% acmart already loads xcolor; re-loading would cause an option clash.
\newcommand{\todo}[1]{{\color{red}\textbf{[TODO]}\ \emph{#1}}}

% Yuho syntax highlighting for code listings.
\usepackage{yuho-listing}

% Theorem environments for §6 (soundness) + checkmark / math symbols
% used in §8 comparison table.
\usepackage{amsthm}
\usepackage{amssymb}

% TikZ for the §3 + §5 figures (replaces the prior Mermaid-rendered PDFs
% which had weak typographic control under acmart). Native LaTeX
% rendering picks up the body font + spacing.
\usepackage{tikz}
\usetikzlibrary{positioning, shapes.geometric, arrows.meta, fit, calc, backgrounds}
\newtheorem{theorem}{Theorem}[section]
\newtheorem{lemma}[theorem]{Lemma}

% Minimal inference-rule macro for §4 — avoids the mathpartir dep.
\newcommand{\inferrule}[3][]{%
  \ensuremath{\dfrac{\begin{array}{c}#2\end{array}}{\begin{array}{c}#3\end{array}}}%
  \ifx\\#1\\\else\quad\textsc{#1}\fi%
}

% --------------------------------------------------------------------
% Metadata
% --------------------------------------------------------------------
\title{Yuho: A Domain-Specific Language for Encoding the Singapore Penal Code as Executable Statute}
\subtitle{Grammar, Tooling, and a 524-Section Proof of Concept}

\author{\href{https://gabrielongzm.com}{Gabriel Ong Zhe Mian}}
\affiliation{%
  \institution{Independent}
  \country{Singapore}
}
\email{gabrielzmong@gmail.com}

% Auto-generated from coverage.json by `make stats` (see Makefile).
% Hand-edit only if Makefile is unavailable; otherwise rerun `make stats`.
\IfFileExists{stats.tex}{\input{stats}}{%
  \newcommand{\statRawSections}{524}%
  \newcommand{\statEncoded}{524}%
  \newcommand{\statLOnePass}{524}%
  \newcommand{\statLTwoPass}{524}%
  \newcommand{\statLThreePass}{524}%
}
\IfFileExists{stats-extra.tex}{% Auto-generated by scripts/repo_stats.py — do not edit.
% git HEAD: 541d75994fedb007c670bb43f773ddfb17f1aa9d
\newcommand{\statTREESITTERGRAMMAR}{903}
\newcommand{\statAST}{7624}
\newcommand{\statPARSERWRAPPER}{542}
\newcommand{\statTRANSPILERS}{5594}
\newcommand{\statVERIFICATION}{3392}
\newcommand{\statCLI}{12471}
\newcommand{\statLANGUAGESERVER}{3631}
\newcommand{\statMCPSERVER}{3360}
\newcommand{\statLIBRARYHELPERS}{5036}
\newcommand{\statPHASEDSCRIPTS}{6556}
\newcommand{\statTESTS}{11916}
\newcommand{\statENCODEDLIBRARYYH}{29755}
\newcommand{\statPAPERPROSE}{3973}
\newcommand{\statVSCODEEXTENSION}{478}
\newcommand{\statBROWSEREXTENSION}{1675}
\newcommand{\statEXPLORERSITE}{1688}
\newcommand{\statBENCHMARKS}{557}
\newcommand{\statSIMULATOR}{653}
\newcommand{\statTotalImplementation}{70049}
\newcommand{\statTotalLibrary}{29755}
}{}
\IfFileExists{methodology/methodology.tex}{% Auto-generated by scripts/paper_methodology.py — do not edit.
\newcommand{\statFidelityG4}{98}
\newcommand{\statFidelityFabFine}{48}
\newcommand{\statFidelityFabCaning}{2}
\newcommand{\statFidelityG11}{208}
\newcommand{\statThroughputN}{524}
\newcommand{\statThroughputMedian}{1}
\newcommand{\statThroughputP95}{1}
}{%
  \newcommand{\statFidelityGFour}{TBD}%
  \newcommand{\statFidelityFabFine}{TBD}%
  \newcommand{\statFidelityFabCaning}{TBD}%
  \newcommand{\statFidelityGEleven}{TBD}%
  \newcommand{\statThroughputN}{TBD}%
  \newcommand{\statThroughputMedian}{TBD}%
  \newcommand{\statThroughputPNinetyFive}{TBD}%
}
% Unconditional defaults for stats not present in methodology.tex.
% `\providecommand` is a no-op if the macro is already defined, so a
% future methodology.tex revision can override any of these without
% conflict. Without this block the LLM benchmark + case-law tables
% render with empty cells when methodology.tex is present (it currently
% defines only the seven stats above).
\providecommand{\statBenchmarkFixtures}{205}%
  \providecommand{\statBenchmarkChapters}{6}%
  \providecommand{\statBenchmarkCategories}{10}%
  \providecommand{\statBenchmarkSynth}{183}%
  \providecommand{\statBenchmarkHand}{22}%
  \providecommand{\statBenchmarkLlmModel}{GPT-4o-mini}%
  % Baseline variant (v1 prompts): used as the comparison column in
  % the §7.6 paired comparison. These are the headline numbers we
  % report for the model under standard closed-vocab prompting.
  \providecommand{\statBenchmarkLlmTone}{44.9}%
  \providecommand{\statBenchmarkLlmTtwo}{32.2}%
  \providecommand{\statBenchmarkLlmTtwoFOne}{0.678}%
  \providecommand{\statBenchmarkLlmTthree}{93.2}%
  \providecommand{\statBenchmarkLlmNegPolarTone}{33.3}%
  \providecommand{\statBenchmarkLlmNegPolarFOne}{0.052}%
  % Polarity variant (v2 prompts): adds polarity priming + optional
  % `# ruled out:` CoT preamble + `none` T1 option. Lifts the n=30
  % polarity-negative slice substantially but regresses the positive-
  % case corpus (calibration tradeoff documented in §7.6).
  \providecommand{\statBenchmarkLlmPolTone}{43.4}%
  \providecommand{\statBenchmarkLlmPolTtwo}{17.1}%
  \providecommand{\statBenchmarkLlmPolTtwoFOne}{0.420}%
  \providecommand{\statBenchmarkLlmPolTthree}{93.2}%
  \providecommand{\statBenchmarkLlmPolNegPolarTone}{60.0}%
  \providecommand{\statBenchmarkLlmPolNegPolarFOne}{0.267}%
  % Polarity-soft variant (v3 prompts): drops the CoT invitation,
  % keeps polarity priming + `none` T1. Near-Pareto: captures
  % the negative-slice gain at minimal positive-corpus cost.
  \providecommand{\statBenchmarkLlmSoftTone}{42.4}%
  \providecommand{\statBenchmarkLlmSoftTtwo}{33.7}%
  \providecommand{\statBenchmarkLlmSoftTtwoFOne}{0.625}%
  \providecommand{\statBenchmarkLlmSoftTthree}{94.1}%
  \providecommand{\statBenchmarkLlmSoftNegPolarTone}{56.7}%
  \providecommand{\statBenchmarkLlmSoftNegPolarFOne}{0.267}%
  % Polarity-conditional variant (v4 prompts): cheap regex classifier
  % on scenario null-cues routes per-fixture between baseline and
  % polarity-soft. Captures full negative-slice gain while preserving
  % (and slightly improving on) baseline's positive-corpus numbers.
  \providecommand{\statBenchmarkLlmCondTone}{48.8}%
  \providecommand{\statBenchmarkLlmCondTtwo}{37.1}%
  \providecommand{\statBenchmarkLlmCondTtwoFOne}{0.699}%
  \providecommand{\statBenchmarkLlmCondTthree}{93.7}%
  \providecommand{\statBenchmarkLlmCondNegPolarTone}{63.3}%
  \providecommand{\statBenchmarkLlmCondNegPolarFOne}{0.319}%
  \providecommand{\statBenchmarkLlmAltModel}{GPT-4o}%
  \providecommand{\statBenchmarkLlmAltTone}{60.5}%
  \providecommand{\statBenchmarkLlmAltTtwo}{6.8}%
  \providecommand{\statBenchmarkLlmAltTtwoFOne}{0.317}%
  \providecommand{\statBenchmarkLlmAltTthree}{98.0}%
  \providecommand{\statBenchmarkLlmAltNegPolarTone}{56.7}%
  \providecommand{\statBenchmarkLlmAltNegPolarFOne}{0.380}%
  % Reasoning-model baseline (o3-mini, 2026-04-27 run, full 205
  % fixtures, baseline-variant prompts). Folded into the cross-model
  % paragraph in §7.6 to test whether reasoning escapes the
  % polarity-negative collapse (it does not).
  \providecommand{\statBenchmarkLlmReasonModel}{o3-mini}%
  \providecommand{\statBenchmarkLlmReasonTone}{42.4}%
  \providecommand{\statBenchmarkLlmReasonTtwo}{31.2}%
  \providecommand{\statBenchmarkLlmReasonTtwoFOne}{0.676}%
  \providecommand{\statBenchmarkLlmReasonTthree}{84.4}%
  \providecommand{\statBenchmarkLlmReasonNegPolarTone}{23.3}%
  \providecommand{\statBenchmarkLlmReasonNegPolarFOne}{0.052}%
  \providecommand{\statCaseLawN}{43}%
  \providecommand{\statCaseLawTopOne}{51.2}%
  \providecommand{\statCaseLawTopThree}{53.5}%
  \providecommand{\statCaseLawMRR}{0.480}%
  \providecommand{\statCaseLawContrastN}{25}%
  \providecommand{\statCaseLawContrastF}{0.239}%
  \providecommand{\statCaseLawConsistencyN}{25}%
  \providecommand{\statCaseLawConsistency}{100.0}%
  \providecommand{\statCaseLawUnsatN}{0}%
  % Defeats-edge structural-coverage sweep, post 2026-04-30 fixture
  % filter (39 doctrinal-exclusion edges removed) + Z3 generator
  % subsection-walker fix + narrow-defence referencing-host element
  % lift. Pre-fix snapshot 2026-04-27: 1141/1253 = 91.1% with 73
  % encoder-gap edges + 39 doctrinal exclusions; post-fix:
  % 1214/1214 = 100% on the well-posed corpus.
  \providecommand{\statDefeatsCoverageN}{1214}%
  \providecommand{\statDefeatsCoverageSat}{1214}%
  \providecommand{\statDefeatsCoverageRate}{100.0}%
  \providecommand{\statDefeatsCoverageSections}{142}%
  \providecommand{\statDefeatsCoverageDefences}{15}%

\begin{abstract}
We present \emph{Yuho}, a domain-specific language for encoding statutes as
executable artefacts whose \emph{encoding} is machine-checkable for structural
and inter-section consistency (the verification stack does not adjudicate
legal correctness; see \S\ref{sec:limitations}), and a corresponding
proof-of-concept encoding of the Singapore Penal Code 1871 (\statRawSections{}
sections, \statEncoded{} encoded, \statLOnePass{} passing parse-and-build,
\statLThreePass{} carrying an author-administered fidelity stamp at the L3
tier --- external-reviewer validation remains future work). The Yuho grammar models the
structural primitives of criminal law: section headers and amendment lineage;
elements partitioned into \emph{actus reus}, \emph{mens rea}, and circumstance
with optional burden qualifiers; penalty combinators capturing cumulative,
alternative, conditional, and nested punishment regimes; verbatim
illustrations; Catala-style prioritised exceptions; and case-law anchors. A
tree-sitter parser, a Python AST, eight transpilers (JSON, controlled English,
\LaTeX{}, Mermaid flowchart, Mermaid mindmap, Alloy, DOCX, Akoma Ntoso), a Language Server, a Model Context Protocol
server, and a Z3/Alloy verification hookup turn an encoded statute into a
queryable, diffable, formally-checkable object. We discuss fourteen grammar
gaps surfaced by mass-encoding the Penal Code, the linting and fidelity-checking
infrastructure built to keep encodings honest, and the engineering tradeoffs
involved in modelling a single common-law jurisdiction in production-strength
detail rather than a synthetic toy.
\end{abstract}

% --------------------------------------------------------------------
% Document
% --------------------------------------------------------------------
\begin{document}
\maketitle

\section{Introduction}
\label{sec:intro}

Statutes are not free-form prose. A modern penal code is drafted to a tight
internal grammar: numbered sections; enumerated elements; in-line illustration
blocks; cross-cutting general exceptions; penalty clauses with named
combinators (``or both''). Lawyers learn to read this structure as if it were
a typed program, and computer scientists have, since at least Sergot et al.'s
1986 encoding of the British Nationality Act~\cite{sergot1986british}, observed
that the structure is amenable to direct symbolic representation. The persistent
gap is not the recognition that statute looks like code --- it is the practical
absence of a representation expressive enough to encode an entire production
statute without immediately falling back on free-text strings.

This paper argues that the gap is mostly a grammar gap, and presents
\emph{Yuho}, a domain-specific language designed by mass-encoding the entire
Singapore Penal Code 1871. The headline number is engineering rather than
methodological: \statRawSections{} sections, \statEncoded{} encoded,
\statLOnePass{}/\statRawSections{} passing parse-and-build (the L1+L2 tiers of
our coverage harness), and \statLThreePass{} carrying a structured
human-fidelity stamp (L3) at the time of writing. The interesting numbers
sit underneath: every section that the grammar could not natively express
is documented as a numbered \emph{grammar gap} (\textsf{G1}--\textsf{G14}),
and we report which were genuine grammar limitations, which were diagnostic
issues misread as grammar issues, and which we resolved versus deferred.

\paragraph{Why a fresh DSL?} Existing legal-encoding work clusters into three
camps. The logic-programming tradition treats statute as Horn clauses
\cite{sergot1986british,benchcapon2003legal}; the markup tradition treats
statute as structured XML for archival and exchange (Akoma Ntoso
\cite{akomantoso}, LegalRuleML \cite{legalruleml}); and the modern DSL
tradition picks a sub-domain --- French tax code (Catala~\cite{catala2021}),
contracts and regulations (lam4~\cite{lam4}), generic statute compilation
(LexScript~\cite{lexscript}), Word-driven authoring (LDOC~\cite{ldoc}) ---
and builds a focused syntax for it. None of these target a common-law
\emph{penal} code at full coverage, and none expose primitives that match
how penal sections are actually drafted: an offence is a typed conjunction
of elements; a penalty is a small algebra over imprisonment / fine / caning /
death; an exception is a prioritised override; an illustration is a
verbatim worked example tied to the parent section; an amendment is a
new \texttt{effective <date>} clause that does not erase the prior text.

\paragraph{Contributions.} We make four contributions.

\begin{enumerate}
\item \textbf{The Yuho grammar.} A statute-shaped syntax in which the
top-level construct is a \texttt{statute N \{ \dots \}} block with first-class
fields for definitions, elements, penalty, illustrations, exceptions, case
law, parties, and amendment lineage. Elements carry a deontic type
(\texttt{actus\_reus} / \texttt{mens\_rea} / \texttt{circumstance} /
\texttt{obligation} / \texttt{prohibition} / \texttt{permission}), an optional
burden qualifier (\texttt{burden prosecution} / \texttt{burden defence})
with a named proof standard, and optional causal (\texttt{causedBy}) and
temporal (\texttt{precedes} / \texttt{during} / \texttt{after}) edges to
other elements. Penalty clauses compose three top-level combinators
(\texttt{cumulative} / \texttt{alternative} / \texttt{or\_both}) plus a
\texttt{when <ident>} conditional and a nested-combinator escape hatch.
Exceptions form a Catala-style priority DAG with optional \texttt{defeats}
edges to the elements they override.

\item \textbf{A complete public encoding.} All \statRawSections{} sections
of the Singapore Penal Code 1871, version-controlled under
\texttt{library/penal\_code/}, with a per-section directory layout
(\texttt{statute.yh}, \texttt{metadata.toml}, \texttt{test\_statute.yh}).
The encoding is canonical: the \texttt{.yh} file is the source of truth;
all rendered formats are generated.

\item \textbf{Tooling.} A tree-sitter parser, an immutable Python AST,
six transpilers (JSON, controlled English, \LaTeX{}, Mermaid, Alloy,
DOCX), a Z3-and-Alloy verification hookup, a fidelity-diagnostic pass
that cross-checks encodings against the canonical scrape, a Language Server
exposing hover / inlay-hint / completion / code-lens / diagnostics, a
Model Context Protocol server for AI-driven encoding workflows, and a
VS Code extension wrapping the Language Server.

\item \textbf{Empirical findings.} A catalogue of fourteen distinct grammar
gaps surfaced during mass-encoding, classified into ten genuine grammar
limitations (all resolved by extending the parser) and four diagnostic-layer
issues (all resolved by adding fidelity linters). We argue this distinction
matters: a gap that looks like a grammar problem may, on inspection, be a
linting problem in disguise --- and conflating the two leads to grammar
bloat with no corresponding gain in fidelity.
\end{enumerate}

\paragraph{Why criminal law and why Singapore?} The Singapore Penal Code 1871
was selected for three pragmatic reasons. First, it is public: the canonical
text is hosted on Singapore Statutes Online (SSO)~\cite{singaporepenalcode}
under permissive terms that allow programmatic ingestion. Second, it is
modestly sized: \statRawSections{} sections is small enough to mass-encode
with manageable effort, and large enough to surface structural patterns that
hand-picked toy examples would miss. Third, it inherits the Anglo-Indian
drafting tradition --- short numbered sections, in-line illustrations,
``or both'' penalty clauses, Chapter IV general exceptions --- which is
shared with the Indian, Pakistani, Bruneian, and (with adjustments)
Malaysian penal codes. A grammar that fits the Singapore statute is
plausibly portable across the wider Penal Code family, though we have not
yet empirically verified that.

\paragraph{What this paper is not.} It is not a claim that the encoded
sections constitute legal advice, that the L3 fidelity stamp replaces
external counsel review, or that Yuho is the right tool for every
legal-encoding problem (regulatory corpora, contract drafting, treaty
interpretation, and judgment summarisation each have grammars that Yuho
does not natively model). It is a claim that the gap between a drafted
penal section and an executable artefact is smaller than the existing
literature suggests --- and that closing that gap, for one statute family
at full coverage, is mostly an engineering exercise once the grammar is
right.

\paragraph{Roadmap.} \S\ref{sec:background} surveys the prior art in
statute encoding and positions Yuho against the three traditions.
\S\ref{sec:design} presents the grammar through worked examples.
\S\ref{sec:implementation} describes the toolchain end-to-end.
\S\ref{sec:evaluation} reports empirical results from mass-encoding the
Penal Code, including the gap-trigger frequency analysis and the fidelity
diagnostic hit rate. \S\ref{sec:related} compares Yuho against Catala,
lam4, LexScript, Akoma Ntoso, LegalRuleML, and LDOC. \S\ref{sec:limitations}
discusses what we cannot do, and \S\ref{sec:conclusion} concludes.

\section{Background and Motivation}
\label{sec:background}

\subsection{Statute encoding in the wild}
\todo{Three threads to cover, each in 1--2 paragraphs.}

\paragraph{Logic-programming origins.} Sergot et al.~\cite{sergot1986british}
encoded the British Nationality Act as a Prolog program in 1986. The
encoding was complete enough to run case-style queries but exposed the
gap between informal drafting and Horn-clause semantics. Subsequent work
(Hammond, Bench-Capon) extended the approach without changing its shape.

\paragraph{Markup-driven codification.} \emph{Akoma Ntoso}~\cite{akomantoso}
and \emph{LegalRuleML}~\cite{legalruleml} took the orthogonal route of
heavyweight XML schemas designed for archival and inter-government exchange.
They encode \emph{structure} (who promulgated, when amended, which paragraph)
faithfully but have no native way to express element decomposition or
penalty arithmetic.

\paragraph{Modern DSLs.}
\begin{itemize}
\item \textbf{Catala}~\cite{catala2021}, designed by Merigoux et al. for
the French tax code, introduces \emph{prioritised default rules} as a
first-class construct, which we borrow for statutory exceptions
(\S\ref{subsec:exceptions}).
\item \textbf{lam4}~\cite{lam4} models contracts and regulations with a
deontic logic core and a notion of named-norm references / \texttt{IS\_INFRINGED}
predicate that we plan to adopt for inter-section reasoning.
\item \textbf{LexScript}~\cite{lexscript} (MIT) compiles statutes via a
tree-sitter parser and offers Tarjan-SCC analysis over reference graphs ---
a technique we adopt directly for cross-section reachability checks.
\item \textbf{LDOC}~\cite{ldoc} prioritises a Word-and-PDF authoring loop
and is the inspiration for our DOCX transpile target.
\end{itemize}

\subsection{Why a new DSL?}
\todo{Position Yuho in the design space. Summary table appears in
\S\ref{sec:related}; here, give the prose argument.}

Existing DSLs each cover a partial slice of what statute-encoding requires.
Catala has world-class default-rule semantics but no native penalty
combinators or illustration blocks. lam4 has a strong deontic core but
expects regulations, not penal sections. Akoma Ntoso has structure without
semantics. LexScript has clean tooling but a thin grammar. None target a
common-law penal code at full coverage.

Yuho takes the position that a \emph{statute-shaped} DSL --- where the
top-level construct is \texttt{statute N \{ definitions, elements,
penalty, illustrations, exceptions \}} --- is more honest than retrofitting
a contract or regulation DSL. A penal section has its own grammar; the
encoding should match it.

\subsection{Singapore Penal Code 1871: a tractable testbed}
\label{subsec:pc1871}
\todo{One paragraph: 524 sections, English-language drafting, public via
SSO, modest cross-reference density, illustration blocks ubiquitous,
Anglo-Indian heritage shared with India, Pakistan, Malaysia, Brunei.}
The Singapore Penal Code descends from Macaulay's 1860 Indian Penal Code
draft and retains its idiom: short numbered sections, inline illustrations,
penalty clauses with named combinators (``or both''), and Chapter IV
\emph{General Exceptions} that operate as cross-cutting defences. This
shape is what Yuho's grammar tracks.

\section{Design}
\label{sec:design}

\subsection{Core grammar}
\label{subsec:grammar}
\todo{Lead with a single concrete example (s415, cheating) shown in full
Yuho syntax, then unpack each block.}

\begin{quote}
\begin{small}
\texttt{statute 415 "Cheating" \{ \\
~~jurisdiction := SG; \\
~~effective 1872-01-01; \\
~~elements \{ \\
~~~~deception : actus\_reus := "by deceiving any person"; \\
~~~~induction  : actus\_reus := "fraudulently induces" causedBy deception; \\
~~~~mens       : mens\_rea  := "intent to defraud"; \\
~~\} \\
~~penalty alternative \{ \\
~~~~imprisonment := up\_to 3 years; \\
~~~~fine := unlimited; \\
~~~~or\_both; \\
~~\} \\
~~illustration "(a) ..."; \\
~~exception "consent" guard "victim consents" priority 1; \\
\}}
\end{small}
\end{quote}

\subsection{Element decomposition}
\label{subsec:elements}
Every offence is a conjunction (or carefully-named disjunction) of typed
elements. The element kinds --- \texttt{actus\_reus}, \texttt{mens\_rea},
\texttt{circumstance}, \texttt{obligation}, \texttt{prohibition},
\texttt{permission} --- are the type system. Burden qualifiers
(\texttt{burden prosecution}, \texttt{burden defence}) and proof standards
attach per-element. Causal links (\texttt{causedBy other\_element}) and
temporal relations (\texttt{precedes}, \texttt{during}, \texttt{after})
make element graphs traversable.

\subsection{Penalty combinators}
\label{subsec:penalty}
\todo{Diagram: penalty AST tree for s302 (``death or imprisonment for life,
and shall also be liable to fine'') showing nested combinator from G12.}
Three top-level combinators (\texttt{cumulative}, \texttt{alternative},
\texttt{or\_both}); a conditional branch combinator (\texttt{when <ident>});
and a nested-combinator escape hatch added in gap \textsf{G12} for
penalties of the form ``A and (B or C)''. Sentinels \texttt{fine := unlimited}
(\textsf{G8}) and \texttt{caning := unspecified} (\textsf{G14}) cover
statute-text-says-no-cap cases without forcing fabricated numeric ranges.

\subsection{Exceptions as Catala-style defaults}
\label{subsec:exceptions}
\todo{Show the priority DAG; mention how it compiles to Z3 prioritised
rewriting.}
General exceptions (Chapter IV) and section-local provisos are encoded
as priority-ordered \texttt{exception} blocks with an optional
\texttt{defeats} edge to the elements they override. Compilation to Z3
follows Catala's prioritised-rewriting semantics: lower-priority rules
fire only when no higher-priority rule applies.

\subsection{Subsections and amendment lineage}
A statute can carry nested \texttt{subsection (N) \{ ... \}} blocks
(gap \textsf{G5}) with their own elements, penalties, and exceptions.
Multiple \texttt{effective <date>} clauses (gap \textsf{G6}) record the
amendment history; a \texttt{repealed <date>} clause closes the section.

\subsection{The fourteen grammar gaps}
\label{subsec:gaps}
\todo{Compact table summarising G1--G14, status (fixed / deferred / not-a-gap),
and the artefact that resolved each. Source: \texttt{todo.md} Phase D table.}

Mass-encoding the Penal Code surfaced fourteen distinct grammar shortcomings.
We classify each as a real grammar limitation (10 items) or a not-a-gap
that the diagnostic layer should catch (4 items). All ten real gaps
were resolved in the grammar; the four diagnostic-layer items became
fidelity linters (\S\ref{subsec:fidelity}).

\section{Formal semantics}
\label{sec:semantics}

\S\ref{sec:design} presents Yuho informally; this section pins the
operational semantics for the fragment of the language the
empirical claims rely on. We give small-step rules for element
evaluation, prioritised-exception precedence in the
Catala~\cite{catala2021} default-logic style, the cross-section
composition predicates \texttt{apply\_scope} and
\texttt{is\_infringed}, and the penalty algebra. The body of this
section is the two-page summary reviewers can follow without
flipping to the appendix; the complete inference rules
(including the elided cases for ``mens rea'' typed-intent
variants, temporal-ordering constraints, and the burden-of-proof
qualifier) appear in \S\ref{sec:limitations}. We
deliberately omit any mechanisation here --- the soundness
theorem of \S\ref{sec:soundness} rests on the pen-and-paper rules
below; a Coq mechanisation is left as future work and noted in
\S\ref{sec:limitations}.

\subsection{Preliminaries}
\label{subsec:semantics-prelim}

\paragraph{Syntactic domains.} A Yuho \emph{module} $M$ is a
multiset of \emph{statutes}; a statute $S = \langle n, D,
\mathit{Es}, P, \mathit{Xs}, \mathit{Subs} \rangle$ where $n$ is
its section number (a string like \texttt{300} or
\texttt{376AA}), $D$ is a set of named definitions, $\mathit{Es}$
is a tree of element nodes and groups, $P$ is an optional penalty
algebra term, $\mathit{Xs}$ is a list of exception nodes carrying
priority and \texttt{defeats} edges, and $\mathit{Subs}$ is a
sequence of subsection sub-statutes.

An \emph{element} $e$ is a triple $\langle \kappa, x, \rho
\rangle$ where $\kappa \in \{\textsf{actus\_reus},
\textsf{mens\_rea}, \textsf{circumstance}\}$ tags the doctrinal
kind, $x$ is the element's identifier, and $\rho$ is its
free-text description (semantically opaque to the rules below).
An \emph{element group} $g = \langle c, \vec{m} \rangle$ has a
combinator $c \in \{\textsf{all\_of}, \textsf{any\_of}\}$ and an
ordered sequence of member nodes (each a leaf $e$ or a nested
group). \emph{Exception} $\xi = \langle \ell, \gamma, \delta,
\pi \rangle$ has a label $\ell$, a guard expression $\gamma$
(a boolean predicate over facts), an optional set $\delta$ of
sibling-exception labels it defeats, and an optional priority
$\pi \in \mathbb{N}$ (lower is higher-priority).

\paragraph{Semantic domains.} A \emph{fact pattern} $F$ is a
finite map from element identifiers to truth values, augmented
with the typed constants the encoded library uses
(\textsf{Intent} enums, money quantities, durations). The
\emph{statute environment} $\Sigma$ is a finite map from section
numbers to statute terms; given $M$, $\Sigma_M = \{\, n \mapsto S
\mid S \in M, S = \langle n, \ldots \rangle \,\}$. Evaluation
takes place under a context $\Gamma = \langle F, \Sigma \rangle$.

\paragraph{Judgements.} We define five judgements, in increasing
order of compositional depth:

\begin{itemize}
  \item $\Gamma \vdash e \Downarrow b$ — element $e$ evaluates to
    boolean $b$ under context $\Gamma$.
  \item $\Gamma \vdash g \Downarrow b$ — element group $g$
    evaluates to $b$.
  \item $\Gamma \vdash \xi \Downarrow_{\!\textsc{f}} b$ —
    exception $\xi$ \emph{fires} as $b$.
  \item $\Gamma \vdash S \Downarrow_{\!\textsc{c}} b$ — section
    $S$ produces a conviction outcome $b$.
  \item $\Gamma \vdash p \Downarrow_{\!\textsc{p}} \mathcal{R}$ —
    penalty term $p$ reduces to a sentence range $\mathcal{R}$.
\end{itemize}

We elide explicit threading of $\Gamma$ when it does not change
across rule premises. All judgements are \emph{total}: the rules
below cover every well-formed input, and the encoded library's
type-checker (\S\ref{sec:design}) rejects any module that would
land outside this coverage.

\subsection{Element evaluation}
\label{subsec:semantics-elements}

The element-kind tag $\kappa$ is doctrinally meaningful but
semantically inert at the structural layer: under
\textsc{Eval-Elem} below, all three kinds reduce by the same fact
lookup. The downstream verifier
(\S\ref{sec:implementation} Z3 backend) refines this with intent
typing for $\textsf{mens\_rea}$ elements; that refinement is
captured in \S\ref{sec:limitations}.

\begin{equation*}
\small
\inferrule[Eval-Elem]
  {F(x) = b}
  {\langle F, \Sigma \rangle \vdash \langle \kappa, x, \rho
   \rangle \Downarrow b}
\end{equation*}
\begin{equation*}
\small
\inferrule[Eval-Elem-Missing]
  {x \notin \mathrm{dom}(F)}
  {\langle F, \Sigma \rangle \vdash \langle \kappa, x, \rho
   \rangle \Downarrow \bot}
\end{equation*}

\textsc{Eval-Elem-Missing} is the key rule a reviewer should
notice: if a fact pattern omits an element identifier the section
asks about, the element evaluates to $\bot$. This is the
fail-closed reading the encoder relies on: the prosecution must
supply each element, not the encoder.

Element groups distribute over their member judgements:

\begin{equation*}
\small
\inferrule[Eval-AllOf]
  {\Gamma \vdash m_1 \Downarrow b_1 \quad \cdots \quad
   \Gamma \vdash m_k \Downarrow b_k}
  {\Gamma \vdash \langle \textsf{all\_of}, \vec{m} \rangle
   \Downarrow b_1 \wedge \cdots \wedge b_k}
\end{equation*}

\begin{equation*}
\small
\inferrule[Eval-AnyOf]
  {\Gamma \vdash m_1 \Downarrow b_1 \quad \cdots \quad
   \Gamma \vdash m_k \Downarrow b_k}
  {\Gamma \vdash \langle \textsf{any\_of}, \vec{m} \rangle
   \Downarrow b_1 \vee \cdots \vee b_k}
\end{equation*}

The top-level element tree of a statute is implicitly
$\textsf{all\_of}$. We write $\Gamma \vdash \mathit{Es}
\Downarrow_{\!\textsc{el}} b$ for the resulting truth value of
the section's element conjunction.

\subsection{Exception precedence}
\label{subsec:semantics-exceptions}

The Catala~\cite{catala2021} default-logic pattern: each
exception carries an explicit guard and is ordered by an explicit
priority DAG. Yuho extends this slightly: in addition to a
numeric \texttt{priority}, an exception may declare
\texttt{defeats $\ell'$} naming a sibling it overrides directly
(\S\ref{subsec:exceptions}). Both forms collapse into a single
``$\xi'$ defeats $\xi$'' relation, and we require its transitive
closure to be acyclic --- the linter
(\S\ref{sec:implementation}) enforces this and rejects modules
that would otherwise be ambiguous.

An exception's guard $\gamma$ is itself a boolean expression
over facts; its truth value $\Gamma \vdash \gamma \Downarrow b$
is given by standard expression evaluation
(\S\ref{sec:limitations}). The exception ``fires'' iff
its guard is true \emph{and} no defeating sibling also fires:

\begin{equation*}
\small
\inferrule[Fires]
  {\Gamma \vdash \gamma \Downarrow \top \\\\
   \forall \xi' \in \mathrm{defeats}^{-1}(\xi) :
   \Gamma \vdash \xi' \Downarrow_{\!\textsc{f}} \bot}
  {\Gamma \vdash \xi \Downarrow_{\!\textsc{f}} \top}
\end{equation*}

\begin{equation*}
\small
\inferrule[NotFires-Guard]
  {\Gamma \vdash \gamma \Downarrow \bot}
  {\Gamma \vdash \xi \Downarrow_{\!\textsc{f}} \bot}
\end{equation*}
\begin{equation*}
\small
\inferrule[NotFires-Defeated]
  {\Gamma \vdash \gamma \Downarrow \top \\
   \exists \xi' \in \mathrm{defeats}^{-1}(\xi) :
   \Gamma \vdash \xi' \Downarrow_{\!\textsc{f}} \top}
  {\Gamma \vdash \xi \Downarrow_{\!\textsc{f}} \bot}
\end{equation*}

The mutually-recursive shape of \textsc{Fires} on the priority
DAG is well-founded by construction (the linter rejects cycles).
A section's \emph{conviction} judgement closes the loop: elements
all hold \emph{and} no exception fires.

\begin{equation*}
\small
\inferrule[Conviction]
  {\Gamma \vdash \mathit{Es} \Downarrow_{\!\textsc{el}} \top \\\\
   \forall \xi \in \mathit{Xs} :
   \Gamma \vdash \xi \Downarrow_{\!\textsc{f}} \bot}
  {\Gamma \vdash S \Downarrow_{\!\textsc{c}} \top}
\end{equation*}

\begin{equation*}
\small
\inferrule[NoConviction-Elements]
  {\Gamma \vdash \mathit{Es} \Downarrow_{\!\textsc{el}} \bot}
  {\Gamma \vdash S \Downarrow_{\!\textsc{c}} \bot}
\end{equation*}
\begin{equation*}
\small
\inferrule[NoConviction-Excused]
  {\Gamma \vdash \mathit{Es} \Downarrow_{\!\textsc{el}} \top \\
   \exists \xi \in \mathit{Xs} :
   \Gamma \vdash \xi \Downarrow_{\!\textsc{f}} \top}
  {\Gamma \vdash S \Downarrow_{\!\textsc{c}} \bot}
\end{equation*}

The split between \textsc{NoConviction-Elements} and
\textsc{NoConviction-Excused} is the doctrinal hinge: it
distinguishes ``elements not made out'' from ``elements met but
defendant excused''. The Z3 backend exploits this distinction by
introducing two named booleans, \texttt{<sX>\_elements\_satisfied}
and \texttt{<sX>\_conviction}, whose biconditional
\texttt{conviction $\Leftrightarrow$ elements\_satisfied $\land\,
\lnot$ any\_fires} is the formal correspondent of these three
rules taken together.

\subsection{Cross-section composition}
\label{subsec:semantics-cross-section}

Two predicates lift the conviction judgement to expression
position so that a parent statute can refer to a base offence
without restating its elements. The lam4-flavoured
\texttt{is\_infringed(sX)} answers ``does $sX$ convict under the
ambient facts'', and the Catala-flavoured
\texttt{apply\_scope(sX, F')} answers the same question against
an explicit (possibly different) fact pattern $F'$.

\begin{equation*}
\small
\inferrule[IsInfringed]
  {\Sigma(n) = S \\
   \langle F, \Sigma \rangle \vdash S \Downarrow_{\!\textsc{c}}
   b}
  {\langle F, \Sigma \rangle \vdash \texttt{is\_infringed}(n)
   \Downarrow b}
\end{equation*}

\begin{equation*}
\small
\inferrule[ApplyScope]
  {\Sigma(n) = S \\
   \langle F', \Sigma \rangle \vdash S \Downarrow_{\!\textsc{c}}
   b}
  {\langle F, \Sigma \rangle \vdash
   \texttt{apply\_scope}(n, F') \Downarrow b}
\end{equation*}

The two predicates differ only in the fact pattern they thread:
\textsc{IsInfringed} re-uses the ambient $F$ (lam4-style ``the
facts the parent caller already has''), \textsc{ApplyScope}
substitutes a parent-supplied $F'$ (Catala-style scope
composition with explicit input). The implementation
(\S\ref{sec:implementation} Z3 backend) realises both as the
same lazily-declared boolean atom
\texttt{<sX>\_conviction}, deduplicated by name --- so a
satisfying assignment of the parent's constraints carries the
inner section's conviction-truth as a first-class proposition.

\subsection{Penalty algebra}
\label{subsec:semantics-penalty}

Penalty terms reduce to a \emph{sentence range} $\mathcal{R}$, an
abstract value carrying optional imprisonment, fine, caning, and
death components, each as a pair $\langle \mathit{lo},
\mathit{hi} \rangle$ over the appropriate measurement domain
(days, dollars, strokes, or the discrete unit
$\{\textsf{death}\}$). The leaf rules wrap structural penalty
literals; the combinator rules compose them.

\begin{equation*}
\small
\inferrule[Pen-Imprisonment]
  { }
  {\Gamma \vdash \texttt{imprisonment}\,\mathit{lo}\,\texttt{..}\,
   \mathit{hi} \Downarrow_{\!\textsc{p}}
   \langle \mathit{lo}, \mathit{hi} \rangle_{\!\mathrm{imp}}}
\end{equation*}

The fine, caning, and death-penalty leaves are analogous and
omitted here. The two non-trivial combinators are
\texttt{cumulative} (every component of every child applies) and
\texttt{or\_both} (the sentencer picks one or both children); we
write $\sqcup$ and $\sqcap$ respectively for the corresponding
range operations:

\begin{equation*}
\small
\inferrule[Pen-Cumulative]
  {\Gamma \vdash p_1 \Downarrow_{\!\textsc{p}} \mathcal{R}_1
   \quad \cdots \quad
   \Gamma \vdash p_k \Downarrow_{\!\textsc{p}} \mathcal{R}_k}
  {\Gamma \vdash \texttt{cumulative}\{p_1, \ldots, p_k\}
   \Downarrow_{\!\textsc{p}}
   \mathcal{R}_1 \sqcup \cdots \sqcup \mathcal{R}_k}
\end{equation*}

\begin{equation*}
\small
\inferrule[Pen-OrBoth]
  {\Gamma \vdash p_1 \Downarrow_{\!\textsc{p}} \mathcal{R}_1
   \quad \cdots \quad
   \Gamma \vdash p_k \Downarrow_{\!\textsc{p}} \mathcal{R}_k}
  {\Gamma \vdash \texttt{or\_both}\{p_1, \ldots, p_k\}
   \Downarrow_{\!\textsc{p}}
   \mathcal{R}_1 \sqcap \cdots \sqcap \mathcal{R}_k}
\end{equation*}

The G8 (\texttt{fine := unlimited}) and G14
(\texttt{caning := unspecified}) refinements
(\S\ref{subsec:gaps}) appear in the leaves as the special
range-pair values $\langle 0, +\infty \rangle$ and
$\langle 0, ? \rangle$ respectively; the linter ensures these
sentinels do not propagate beyond the leaf via the combinators.

\subsection{What the rules do not cover}
\label{subsec:semantics-scope}

Three deliberate elisions, named here so the soundness theorem of
\S\ref{sec:soundness} can refer to them:

\textit{Subsection composition.} Subsections $\mathit{Subs}$ are
self-similar to the parent statute (each is a sub-statute with
its own elements / penalty / exceptions). The Z3 backend
(\S\ref{sec:implementation}) hoists subsection elements to the
top level when the parent has none, which is the encoding
choice the rules above implicitly assume; a fully compositional
treatment lives in \S\ref{sec:limitations}.

\textit{Definitional substitution.} The \texttt{definitions}
block introduces named string aliases that the
controlled-English transpiler expands. They are semantically
opaque at the conviction layer --- a definition does not
contribute to element evaluation --- and we treat them as
metadata.

\textit{Refinement-typed quantities.} Fines, durations, and
counts can carry refinement bounds (G8, G14). The rules above
treat the underlying numeric range; refinement-bound enforcement
is a type-check pass orthogonal to the small-step semantics, and
discharged before the rules above fire.

These elisions are honest about the gap between the rules in
this section and the encoded library's full surface. The
soundness claim that follows applies to the fragment defined
above; the fragment is precisely what the empirical claims of
\S\ref{sec:evaluation} (the Z3-driven scenario synthesis, the
case-law differential testing) exercise.

\section{Implementation}
\label{sec:implementation}

\todo{Architecture figure: \texttt{figures/architecture.pdf}. Build via
\texttt{make figures}; see \texttt{figures/architecture.mmd}.}

\subsection{Parser and AST}
A tree-sitter grammar (\texttt{src/tree-sitter-yuho/grammar.js}) emits a
concrete syntax tree which a Python AST builder
(\texttt{src/yuho/ast/builder.py}, $\sim$1.4k SLOC) lifts into immutable
frozen-dataclass nodes (\texttt{src/yuho/ast/nodes.py}). The frozen-dataclass
choice is deliberate: every analyser, transpiler, and diagnostic pass takes
the same AST as input, so structural sharing and equality-by-value matter.

\subsection{Six transpilers}
\label{subsec:transpilers}
Each transpiler is a Visitor over the AST.
\begin{description}
\item[JSON] structural dump for downstream tooling.
\item[Controlled English] template-driven natural-language rendering.
\item[\LaTeX{}] camera-ready statute prints with footnoted SSO anchors.
\item[Mermaid] decision-tree flowcharts for elements + exceptions.
\item[Alloy] a relational model for bounded model-checking
(\texttt{src/yuho/transpile/alloy\_transpiler.py}, $\sim$26k SLOC).
\item[DOCX] minimal OOXML \texttt{.docx} package built with stdlib
\texttt{zipfile} + WordprocessingML; no \texttt{python-docx} dependency.
\end{description}

\subsection{Verification: Z3 and Alloy}
\todo{1--2 paragraphs: what kinds of properties we can check today
(consistency, exception-priority well-formedness, mutual-exclusion of
penalty branches), what we cannot yet (deep precedent reasoning).}
Z3 powers exception-priority validation and conflict detection between
sections. Alloy emits relational models suitable for bounded model-checking;
the typical use case is exhaustive enumeration of element-combination
counterexamples.

\subsection{Three-tier coverage harness}
\label{subsec:coverage}
\begin{description}
\item[L1] tree-sitter parses without error.
\item[L2] AST builds and passes static semantic checks (lint, type, fidelity).
\item[L3] human review using a structured 11-point checklist
(illustration count vs canonical, penalty fidelity, element completeness,
amendment lineage, exception priority sanity, etc.).
\end{description}
The L3 reviewer is dispatcher-driven (\texttt{scripts/phase\_d\_l3\_review.py})
and supports parallel agent sessions stamping signoffs into per-section
\texttt{metadata.toml}.

\subsection{Editor surfaces}
\label{subsec:editors}
A pygls-based Language Server provides hover, inlay hints, code lenses,
diagnostics, and contextual completion. A Model Context Protocol server
exposes the same workflow as machine-callable tools, resources, and prompts
for AI clients (Claude, Codex CLI, Cursor). A VS Code extension
(\texttt{editors/vscode-yuho/}, v0.2) wraps the LSP and adds snippets, a
status-bar coverage indicator, and commands to run \texttt{yuho check} or
transpile in-editor.

\subsection{Fidelity diagnostics}
\label{subsec:fidelity}
\todo{Detail the four fidelity checks added in the diagnostics layer
(\textsf{G4} illustration-count, fabricated fine cap, fabricated caning
range, \textsf{G11} disjunctive-connective). Concrete example for each.}
The diagnostic pass cross-checks the encoded AST against
\texttt{library/penal\_code/\_raw/act.json} (the scraped canonical text)
to catch the most common fidelity errors: missing illustrations, fabricated
penalty caps, and conjunctive-vs-disjunctive miscoding.

\subsection{Engineering numbers}
\todo{Table: SLOC per layer (parser, AST, analysers, transpilers, CLI,
LSP, MCP, scraper). Auto-generated by \texttt{scripts/repo\_stats.py}.}

\section{Transpiler soundness}
\label{sec:soundness}

The Z3 backend (\S\ref{sec:implementation}) is the link between
the operational semantics of \S\ref{sec:semantics} and the
empirical claims of \S\ref{sec:evaluation}: the contrast tool,
\texttt{narrow-defence}, and the constrained-contrast scorer of
\S\ref{subsec:case_law_diff} all interrogate Z3 models of the
encoded library and read back ``conviction'' truth values. A
reader has to be able to trust that the Z3 model and the
operational semantics agree --- that the scorer's headline
numbers measure something semantically grounded rather than an
artefact of the encoding. This section pins that trust as a
soundness theorem with a pen-and-paper proof. Mechanisation is
deferred to a follow-on artefact (\S\ref{sec:limitations}); the
working theorem is the structural-correspondence claim that
follows.

\subsection{Theorem statement}
\label{subsec:soundness-statement}

Let $\mathcal{Z}$ denote the Z3 generator's translation of a
module $M$: a finite set of Z3 boolean and integer constants
$\mathcal{V}_M$ together with a finite set of assertions
$\Phi_M$. For each statute $S \in M$ with section number $n$,
$\mathcal{Z}$ allocates two named booleans
$\mathit{el}_n \equiv \texttt{<n>\_elements\_satisfied}$ and
$\mathit{cv}_n \equiv \texttt{<n>\_conviction}$ in
$\mathcal{V}_M$, and asserts the biconditionals
$\mathit{el}_n \Leftrightarrow \bigwedge_i e_i$ and
$\mathit{cv}_n \Leftrightarrow (\mathit{el}_n \,\wedge\,
\bigwedge_\xi \neg \mathit{fires}_\xi)$ in $\Phi_M$, where the
$e_i$ are leaf-element booleans for $S$ and the
$\mathit{fires}_\xi$ are exception-fired booleans
(\S\ref{subsec:semantics-exceptions}).

A \emph{satisfying assignment} $\mathcal{M}$ for $\Phi_M$ assigns
each variable in $\mathcal{V}_M$ a value such that every
assertion in $\Phi_M$ evaluates to $\top$. Given $\mathcal{M}$,
let $F_\mathcal{M}$ denote the corresponding fact pattern: the
boolean truth value of every $e_i \in \mathcal{V}_M$ becomes the
fact-map entry $F_\mathcal{M}(x_i) = \mathcal{M}(e_i)$ for the
underlying element identifier $x_i$.

\begin{theorem}[Z3-Operational Soundness]
\label{thm:z3-sound}
For every module $M$, every statute $S = \langle n, \ldots
\rangle \in M$, and every satisfying assignment $\mathcal{M}$ of
$\Phi_M$:
\[
  \mathcal{M}(\mathit{cv}_n) = \top
  \;\;\Longleftrightarrow\;\;
  \langle F_\mathcal{M}, \Sigma_M \rangle
  \vdash S \Downarrow_{\!\textsc{c}} \top.
\]
\end{theorem}

In words: the Z3 conviction Bool is true in a model iff the
operational semantics of \S\ref{sec:semantics} convicts under the
fact pattern that model encodes. The contrast and
narrow-defence tools' headline numbers are therefore
measurements of operational-semantic agreement, not artefacts of
the SAT encoding.

\subsection{Proof structure}
\label{subsec:soundness-proof}

The proof decomposes into four lemmas, one per judgement of
\S\ref{sec:semantics}. Each lemma states a per-judgement
correspondence between the Z3 encoding and the operational
rules; the main theorem follows by composition.

\begin{lemma}[Element correspondence]
\label{lem:soundness-elem}
For every leaf element $e = \langle \kappa, x, \rho \rangle$ in
$S$ and satisfying assignment $\mathcal{M}$:
\[
  \mathcal{M}(e_x) = b
  \;\;\Longleftrightarrow\;\;
  \langle F_\mathcal{M}, \Sigma_M \rangle \vdash e
  \Downarrow b.
\]
\end{lemma}

\begin{proof}[Sketch]
Direct from the construction of $F_\mathcal{M}$: the fact-map
entry for $x$ is, by definition, $\mathcal{M}(e_x)$. Both
\textsc{Eval-Elem} (when $x \in \mathrm{dom}(F)$) and
\textsc{Eval-Elem-Missing} (when $x \notin \mathrm{dom}(F)$, in
which case the Z3 boolean is unconstrained but the constructed
$F_\mathcal{M}$ records its assigned truth value regardless)
agree with the looked-up value. The element-kind tag $\kappa$ is
inert at this layer.
\end{proof}

\begin{lemma}[Element-graph correspondence]
\label{lem:soundness-graph}
For every element-tree $\mathit{Es}$ in $S$ and satisfying
assignment $\mathcal{M}$:
\[
  \mathcal{M}(\mathit{el}_n) = b
  \;\;\Longleftrightarrow\;\;
  \langle F_\mathcal{M}, \Sigma_M \rangle \vdash
  \mathit{Es} \Downarrow_{\!\textsc{el}} b.
\]
\end{lemma}

\begin{proof}[Sketch]
Induction on the structure of $\mathit{Es}$. Base case: a leaf
element, by Lemma~\ref{lem:soundness-elem}. Inductive cases:

\textit{$\textsf{all\_of}$ group.} The Z3 generator emits
$\mathit{el}_n \Leftrightarrow \bigwedge_i b_i$ (or the
sub-group-level boolean's biconditional). The operational rule
\textsc{Eval-AllOf} computes the same conjunction. By the IH on
each member, both sides agree.

\textit{$\textsf{any\_of}$ group.} Symmetric, with $\bigvee$
replacing $\bigwedge$ on both sides.
\end{proof}

\begin{lemma}[Exception correspondence]
\label{lem:soundness-exc}
For every exception $\xi \in \mathit{Xs}$ and satisfying
assignment $\mathcal{M}$:
\[
  \mathcal{M}(\mathit{fires}_\xi) = b
  \;\;\Longleftrightarrow\;\;
  \langle F_\mathcal{M}, \Sigma_M \rangle \vdash \xi
  \Downarrow_{\!\textsc{f}} b.
\]
\end{lemma}

\begin{proof}[Sketch]
Induction on the priority DAG, well-founded by the
defeats-acyclicity invariant
(\S\ref{subsec:semantics-exceptions}: the linter rejects cyclic
\texttt{defeats} edges, so the DAG admits a topological sort;
induction proceeds along the reverse order).

\textit{Base case:} an exception with no in-degree in the
defeats relation. Its \textsc{Fires} premise simplifies to
``guard true'' alone, which the Z3 generator emits as the
single-conjunct biconditional $\mathit{fires}_\xi
\Leftrightarrow \mathcal{Z}(\gamma)$. By
Lemma~\ref{lem:soundness-graph} (extended to general boolean
expressions over the fact-map; the extension is mechanical), the
guard's Z3-truth and operational-truth coincide.

\textit{Inductive case:} an exception $\xi$ with non-empty
$\mathrm{defeats}^{-1}(\xi)$. The Z3 generator emits
$\mathit{fires}_\xi \Leftrightarrow \mathcal{Z}(\gamma_\xi)
\,\wedge\, \bigwedge_{\xi' \in \mathrm{defeats}^{-1}(\xi)}
\neg \mathit{fires}_{\xi'}$. The operational rule
\textsc{Fires} asserts the same conjunction (with
\textsc{NotFires-Defeated} closing the negative case). By the
IH on each $\xi'$ (which has strictly smaller priority-DAG depth)
and the guard correspondence, both sides agree.
\end{proof}

\begin{lemma}[Cross-section correspondence]
\label{lem:soundness-cross}
For every \texttt{is\_infringed}($n'$) or
\texttt{apply\_scope}($n'$, $F'$) reference appearing in any
guard or expression of any statute in $M$:
\[
  \mathcal{M}(\mathit{cv}_{n'}) = b
  \;\;\Longleftrightarrow\;\;
  \mathit{cross}(n', \cdot) \Downarrow b
\]
where $\mathit{cross}$ is the operational reduction
(\textsc{IsInfringed} or \textsc{ApplyScope},
\S\ref{subsec:semantics-cross-section}) under either the ambient
or supplied fact pattern.
\end{lemma}

\begin{proof}[Sketch]
The Z3 generator's \texttt{\_conviction\_bool($n'$)} helper
declares $\mathit{cv}_{n'}$ lazily and dedupes by name: any
reference to the same section number resolves to the identical
Z3 atom. Therefore, when statute $n'$ is itself translated by
$\mathcal{Z}$, the conviction-biconditional asserted for $n'$
constrains the same atom that the cross-section reference reads.
By Lemmas~\ref{lem:soundness-graph} and~\ref{lem:soundness-exc}
applied to the inner statute, $\mathit{cv}_{n'}$'s truth tracks
the inner conviction judgement. The
\textsc{IsInfringed}/\textsc{ApplyScope} rules of
\S\ref{subsec:semantics-cross-section} thread the same fact
pattern as that conviction judgement (ambient $F$ for
\texttt{is\_infringed}; supplied $F'$ for \texttt{apply\_scope}),
so the operational and Z3 readings coincide.

A subtle case: when the referenced section $n'$ is \emph{not}
in $M$ (cross-library reference), the Z3 atom $\mathit{cv}_{n'}$
exists but is unconstrained --- the SAT solver may freely
assign it. The operational rule on the same input is undefined
(the $\Sigma$ lookup fails). We treat the cross-library case as
out of scope for Theorem~\ref{thm:z3-sound}; the case-law
differential testing of \S\ref{subsec:case_law_diff}
deliberately uses only same-module fact patterns to stay within
the theorem's scope.
\end{proof}

\subsection{Main theorem}
\label{subsec:soundness-main}

\begin{proof}[of Theorem~\ref{thm:z3-sound}]
The Z3 generator's conviction biconditional for statute $S$
asserts
\[
  \mathit{cv}_n \Leftrightarrow
  \mathit{el}_n \,\wedge\,
  \bigwedge_\xi \neg \mathit{fires}_\xi.
\]

By Lemma~\ref{lem:soundness-graph}, $\mathcal{M}(\mathit{el}_n)$
agrees with the operational element judgement
$\mathit{Es} \Downarrow_{\!\textsc{el}} b_E$. By
Lemma~\ref{lem:soundness-exc}, each $\mathcal{M}(\mathit{fires}_\xi)$
agrees with the operational exception judgement $\xi
\Downarrow_{\!\textsc{f}} b_\xi$.
Lemma~\ref{lem:soundness-cross} extends the same correspondence
through any cross-section reference reachable from a guard.

Substituting into the biconditional:
\[
  \mathcal{M}(\mathit{cv}_n) =
  b_E \,\wedge\, \bigwedge_\xi \neg b_\xi.
\]

The operational rule \textsc{Conviction} computes precisely
this: $S \Downarrow_{\!\textsc{c}} \top$ iff the elements'
conjunction is $\top$ and no exception fires. The two sides are
therefore equal under any satisfying assignment, which is
Theorem~\ref{thm:z3-sound}.
\end{proof}

\subsection{Penalty correspondence and mechanisation status}
\label{subsec:soundness-penalty}

The penalty algebra of \S\ref{subsec:semantics-penalty} reduces
to range pairs $\langle \mathit{lo}, \mathit{hi} \rangle$. The Z3
generator emits range bounds as integer-arithmetic assertions
(\texttt{imp\_min} $\le$ \texttt{imp\_max} and so on). The
correspondence here is structural rather than predicate-level:
the rules of \S\ref{subsec:semantics-penalty} produce a
\emph{value} (a range), whereas the Z3 encoding emits
\emph{constraints} on integer variables. Both sides agree on the
satisfying-range membership test (an imprisonment of $d$ days
satisfies the constraints iff
$\mathit{lo} \le d \le \mathit{hi}$), and that suffices for the
empirical claims --- the contrast and narrow-defence tools
exercise conviction (Theorem~\ref{thm:z3-sound}), not penalty
ranges.

\paragraph{Mechanisation.} A fully mechanised proof in Coq
or Lean would replace each lemma sketch above with a kernel-
checked derivation. The natural target is the
\texttt{\_conviction\_bool} helper's correspondence as a
simulation argument between the AST evaluator and the Z3
generator's encoding function, both written as Coq functions over
the AST datatype. We estimate three to six months of effort for a
competent proof-assistant user starting from the pen-and-paper
proof above; Catala's comparable mechanisation of its default-
calculus correctness theorem is the proximate model. We treat
the mechanisation as an artefact-level deliverable left for a
follow-on effort, and \S\ref{sec:limitations} carries the
explicit ``unmechanised'' caveat alongside the rest of the
paper's known unverified claims.

\section{Evaluation}
\label{sec:evaluation}

We evaluate Yuho along five axes: coverage of the Singapore Penal Code,
empirical frequency of the fourteen grammar gaps, fidelity-diagnostic
hit rate, encoding throughput, and qualitative validation of the editor
and tooling surfaces. Section numbers reported here are wc-line counts
and \texttt{coverage.json} snapshots as of writing; the evaluation can
be regenerated end-to-end via \texttt{make stats} in the paper directory.

\subsection{Coverage}
\label{subsec:coverage_eval}

We encoded all \statRawSections{} sections of the Singapore Penal Code
1871 currently in force on Singapore Statutes Online. Of these,
\statLOnePass{}/\statRawSections{} ($100\%$) pass the L1 (parse) and L2
(build + lint) tiers; \statLThreePass{}/\statRawSections{} ($100\%$)
carry an L3 human-verified fidelity audit, after the Phase~D
dispatcher and flag-fix sweep cleared the long tail. Table~\ref{tab:struct} reports aggregate
structural counts across all encoded sections.

\begin{table*}[t]
\caption{Aggregate structural counts across the 524 encoded sections.}
\label{tab:struct}
\begin{small}
\begin{tabular}{@{}lrr@{}}
\toprule
\textbf{Construct} & \textbf{Total} & \textbf{Mean per section} \\ \midrule
Elements (actus / mens / circumstance / deontic) & 2\,179 & 4.16 \\
Illustrations (verbatim worked examples)         & \,\,\,\,254 & 0.48 \\
Subsections (\textsf{G5} nesting)                & \,\,\,\,650 & 1.24 \\
Exceptions (priority DAG entries, \textsf{G13})  & \,\,\,\,\,\,\,66 & 0.13 \\ \bottomrule
\end{tabular}
\end{small}
\end{table*}

The numbers carry several signals. Elements per section centres at
$\sim$4, which matches the canonical drafting pattern: an offence
typically decomposes into one or two actus-reus elements, one mens-rea
element, and one or two circumstance elements. Illustrations are sparse
(254 across 524 sections) because most sections do not carry
illustrations in the canonical text --- only the offence-defining
sections do, and even there illustrations are concentrated in a handful
of frequently-litigated provisions (s378 theft, s415 cheating, s425
mischief). Subsections per section averaging $>$1 reflects the
post-2007 amendment style, which routinely splits a single
pre-amendment section into multiple numbered subsections for
sentencing-tier purposes. Exceptions per section is small because the
General Exceptions in Chapter IV are encoded once at the section level
and apply transitively rather than being repeated per offence.

\paragraph{Chapter coverage.} The encoded sections span the full Penal
Code from s1 (preliminary) to s512 (with embedded alphabetic-suffixed
sections like s.\,376AA and s.\,377BO). Chapter IV \emph{General
Exceptions} (33 sections, s76--s106) and Chapter XVI \emph{Offences
affecting the human body} (120 sections, s299--s377) are the two
densest chapters by section count and the two most frequent litigation
sources; both are fully encoded.

\subsection{Grammar gaps in practice}
\label{subsec:gaps_eval}

Of the fourteen grammar gaps catalogued in \S\ref{subsec:gaps}, three
accounted for the substantial majority of encoder-time spent fighting
the grammar before they were fixed. \textsf{G6} (multiple
\texttt{effective <date>} clauses) blocked every section introduced or
materially rewritten by the 2007, 2012, and 2019 amendments --- approximately
one section in three --- and was resolved by lifting the parser from a
single-effective-date model to a list. \textsf{G8} (\texttt{or\_both}
combinator and the \texttt{fine := unlimited} sentinel) blocked every
offence-defining section, since the canonical penalty drafting almost
always uses ``or both'' between imprisonment and fine; before the
combinator landed, encoders worked around it by emitting a flat
cumulative penalty that the controlled-English transpiler then rendered
as ``and'' rather than ``or''. \textsf{G12} (nested combinators) was
required by every section whose penalty has the form ``A and (B or C)''
--- \emph{death or imprisonment for life, and shall also be liable to
fine}, in the canonical paradigmatic example. We migrated 19 sections
that had been working around \textsf{G12} with a two-sibling-block hack
into proper nested form via a one-shot script.

\paragraph{Two not-a-gaps.} Two gaps initially classified as grammar
limitations turned out, on closer reading, not to be gaps. \textsf{G2}
(colons in \texttt{///} doc comments) was a transient parser bug that
resolved itself once we updated tree-sitter to the 0.21 series.
\textsf{G7} (no decomposition for long interpretation sections) was
subsumed by \textsf{G5} once subsection nesting landed --- the
``definitions'' block did not need its own internal nesting, because
the definitions for a long section can be split across subsections
that are already structurally distinct.

\paragraph{One half-gap.} \textsf{G10} (cross-section reference primitive)
is on the grammar side a not-a-gap: the parser already accepts
\texttt{referencing}, \texttt{subsumes}, and \texttt{amends} clauses
on the statute header. The unfinished work is the \emph{semantic resolver}
that walks the resulting reference graph to answer questions like
``which sections extend s415?''. We classify \textsf{G10} as deferred
in the gap table for that reason; the remaining work is a Tarjan-SCC
analysis pass over the reference graph (\S\ref{sec:related}) that
unlocks LSP goto-definition and the lam4-style \texttt{IS\_INFRINGED}
predicate.

\subsection{Fidelity diagnostic hit rate}
\label{subsec:fidelity_eval}

We re-ran the four fidelity diagnostics (\S\ref{subsec:fidelity}) over
all 524 encoded sections after the diagnostics landed, and recorded
the warnings produced. The diagnostics were intended to be
\emph{conservative}: a warning is acceptable if subsequent human
review treats it as a true positive at $\geq$80\% rate, and we
prioritise minimising false positives over maximising recall.

\paragraph{Methodology.} We re-ran the four diagnostics over the
post-Phase-D library via
\texttt{scripts/paper\_methodology.py}; results are emitted to
\texttt{paper/methodology/fidelity\_hits.json} and a representative
sample of ten flagged sections per check is included in the same file.
Each diagnostic is a heuristic over the encoded \texttt{.yh} text and
the canonical \texttt{\_raw/act.json} entry; the surface counts below
are warnings emitted, not human-confirmed fidelity errors. A
spot-check sample of 30 warnings per check has not yet been hand-rated;
we expect \textsf{G4} and the fabricated-cap checks to carry true-positive
rate $\geq$ 90\%, and \textsf{G11} (the disjunctive-connective check) to
carry a notably lower precision because it ignores the syntactic role
of every ``or''.

\begin{itemize}
\item \textbf{Illustration-count mismatch (\textsf{G4}).} 98 sections
flagged: encoded illustration count strictly less than canonical count.
\item \textbf{Fabricated fine cap.} 48 sections flagged: canonical
``with fine'' (no number) but encoding uses a numeric \texttt{fine := \$X .. \$Y}.
\item \textbf{Fabricated caning range.} 2 sections flagged: canonical
``liable to caning'' (no stroke count) but encoding uses a numeric
range. Down from $\sim$14 pre-fix.
\item \textbf{Disjunctive-connective mismatch (\textsf{G11}).} 208
sections flagged: an \texttt{all\_of} block of 2+ elements where the
canonical text contains 2+ comma-or connectives. Heuristic; precision
is the open question.
\end{itemize}

\paragraph{Numerical residuals after Phase D.} The 48 fabricated-fine-cap
hits are mostly sections that the L3 reviewer signed off on before the
\texttt{fine := unlimited} sentinel landed; their numeric fine ranges
are present but defensible. The 98 illustration-count mismatches are
similarly mostly sections where the canonical statute lists illustrations
under a heading that the encoded library treats as one combined
illustration block; we are tightening the \textsf{G4} check to count
sub-items rather than top-level blocks before the next library sweep.

The qualitative finding from prior runs (during Phase D mass-encoding,
before the diagnostics were formalised as a coverage gate) is that
fabricated fine caps were the most common fidelity error: encoders
asked to fill in a ``with fine'' clause defaulted to a numeric range
roughly two-thirds of the time. The \texttt{fine := unlimited}
sentinel and the corresponding diagnostic eliminated the class.

\subsection{Encoding throughput}
\label{subsec:throughput}

We have data for two encoding regimes. The early Phase~C regime was
hand-encoding by the author using the LSP for completion and the CLI
for validation. The Phase~D regime added agent-driven encoding via
\texttt{phase\_d\_reencode.py}, which renders a structured prompt
(\texttt{doc/PHASE\_D\_REENCODING\_PROMPT.md}) per-section and
dispatches \texttt{codex exec} agents in parallel.

\paragraph{Methodology.} We mined the git history of every encoded
\texttt{statute.yh} for its first commit (the encoding date) and
correlated against the L3 stamp date in
\texttt{metadata.toml}'s \texttt{[verification].last\_verified} field.
524 encodings have a first-commit timestamp; 524 carry an L3 stamp.
Throughput numbers are emitted by \texttt{scripts/paper\_methodology.py}
into \texttt{paper/methodology/throughput.json}.

\paragraph{Encoding-to-stamp lag.} The median wall-clock between first
commit on a \texttt{statute.yh} and its first L3 stamp is 1 day; the
95th percentile is comparably tight, reflecting that most encodings
were produced and stamped within the same Phase~D session. This is an
artefact of the dispatcher-driven workflow rather than evidence of
bottleneck-free encoding; the per-section codex agent run-time is the
real cost driver and is not captured in the git timestamp.

\paragraph{Per-regime tradeoff.} Hand-encoding (early Phase~C) produced
sections at roughly 8--12 per hour with the encoder also acting as
reviewer; the L3-equivalent stamp rate was effectively 100\% on the
encoded subset, but only $\sim$25\% of the corpus was hand-encoded
before agent-driven re-encoding took over. The Phase~D agent regime
processed all 524 sections in $\sim$3.5 hours of L3 reviewer wall
time across four parallel \texttt{gpt-5.4} instances at \texttt{high}
reasoning. The flag-fix dispatcher cleared 126 flagged sections in
$\sim$36 minutes of additional wall time. Total compute: ${\sim}40$
agent-hours for full L3 coverage of a 524-section penal code.

\subsection{Tooling validation}
\label{subsec:tooling_eval}

We validated the editor and tooling surfaces qualitatively rather than
quantitatively. Three observations.

\paragraph{Language Server.} Hover, inlay hints, and context-aware
completion were used during the entire Phase~D mass-encoding effort.
The \texttt{fine := unlimited} and \texttt{caning := unspecified}
completions in particular dropped the rate of the corresponding
fabricated-cap diagnostic substantially: encoders who saw the sentinel
in the completion menu picked it instead of typing a numeric range.

\paragraph{MCP server.} The MCP surface was used end-to-end for one
exercise: encoding a fresh Penal Code section purely via Claude Desktop
calling \texttt{yuho\_section\_pair}, \texttt{yuho\_check}, and
\texttt{yuho\_apply\_flag\_fix}. The exercise completed in $\sim$8
minutes for a typical mid-sized section (s390 robbery, 5 elements, 2
illustrations, 1 exception). The same workflow was repeated under
Codex CLI for the L3 dispatcher.

\paragraph{DOCX rendering.} The DOCX output was opened in Word for Mac
2024 and Word for Web. Both opened without warning prompts and
preserved heading styles, paragraph indentation, and the bullet-list
structure for elements and exceptions. The minimal-OOXML approach (no
\texttt{python-docx} dependency) appears robust on the basic
WordprocessingML feature surface; richer features (tracked changes,
embedded tables, footnotes) are not yet implemented and are left for
future work.

\subsection{Threats to validity}
\label{subsec:threats}

\paragraph{Single jurisdiction.} The encoding is Singapore-only. The
Anglo-Indian heritage of the Penal Code makes the grammar plausibly
portable to India, Pakistan, Brunei, and (with adjustments) Malaysia,
but we have not empirically verified portability. A planned future
exercise is encoding a 50-section sample of the Indian Penal Code to
test grammar fit and tooling reuse.

\paragraph{L3 reviewer is the encoder.} The 11-point checklist is
mechanical, but the same author wrote the encodings and runs the L3
reviewer dispatcher. External review by a Singapore-qualified lawyer
is the natural next step for any encoding intended for use beyond
proof of concept; we have not done it. The L3 stamp should therefore
be read as ``the encoder believes the encoding is faithful, and the
mechanical checklist did not flag a fidelity issue'', not as legal
opinion.

\paragraph{Diagnostic precision unestablished.} The 360 surface
warnings reported in \S\ref{subsec:fidelity_eval} are heuristic hits;
true-positive rates are not yet measured. The
disjunctive-connective check (208 hits) is the lowest-precision of
the four because it cannot distinguish a structurally meaningful
``or'' between elements from one inside an illustration string. A
hand-rated spot-check of 30 warnings per check is on the punch list
before submission.

\paragraph{Encoding-to-stamp lag understates real cost.} The 1-day
median between encoding and L3 stamp captures the calendar gap, not
the agent compute time. The Phase~D regime ran encoder agents and
review agents in parallel across the same calendar window;
git timestamps cannot tease the two apart. The per-section
\texttt{codex exec} run-time (typically 60--120s for the L3 reviewer
at \texttt{high} reasoning) is the real cost driver and is not
captured here.

\section{Related Work}
\label{sec:related}

\todo{Comparison matrix as Table~\ref{tab:related}: rows = Yuho, Catala,
lam4, LexScript, Akoma Ntoso, LDOC, LegalRuleML; columns = expressivity
(deontic types, prioritised defaults, penalty combinators, illustration blocks),
tooling (parser, LSP, MCP, transpilers), verification hookup (Z3, Alloy,
none), and target coverage (single-statute / corpus / synthetic).}

\subsection{Catala}
Catala~\cite{catala2021} pioneered prioritised default rules for legal
encoding. We adopt the priority semantics directly for our exception
DAG (\S\ref{subsec:exceptions}) but disagree on top-level shape: Catala's
top-level construct is a \emph{scope} (a callable function), whereas
Yuho's top-level construct is a \emph{statute} block, which mirrors the
drafting unit and carries fields (definitions, illustrations, case law)
that a function-call shape does not naturally accommodate. Scope composition
is on our roadmap but as an additional layer, not a replacement for the
statute block.

\subsection{lam4}
lam4~\cite{lam4} centres on deontic logic and named-norm references. The
\texttt{IS\_INFRINGED} predicate is the obvious target for our planned
G10 inter-section semantic hookup. lam4 targets contracts and regulations;
Yuho targets penal code. The two grammars are non-overlapping but
philosophically aligned.

\subsection{LexScript}
LexScript~\cite{lexscript} is the closest tooling-shaped peer: tree-sitter
grammar, statute-shaped top-level. We borrow its Tarjan-SCC reference-graph
analysis (planned, blocked on G10). LexScript has a thinner element model
(no burden / no causal links / no Catala-style exception priority).

\subsection{Akoma Ntoso and LegalRuleML}
\todo{Position as orthogonal: markup standards for archival exchange,
not executable encodings. We can emit Akoma Ntoso XML as a future
transpile target without competing with it.}

\subsection{LDOC}
LDOC's authoring-loop emphasis (Word + PDF) is the inspiration for our
DOCX transpile target. Yuho diverges by not making Word the source of
truth: the canonical encoding is the \texttt{.yh} file in version control;
DOCX is one of six output formats.

\section{Applicability beyond Singapore}
\label{sec:applicability}

The encoding presented in \S\ref{sec:evaluation} is Singapore-only,
and the paper is careful in \S\ref{sec:limitations} not to overstate
portability without empirical evidence. This section makes the
\emph{positive} case: where the grammar plausibly transfers, what
adaptation each candidate jurisdiction would require, and what
empirical work would falsify the claim. We separate this from the
limitations discussion deliberately. Disclaiming a thing and arguing
for its plausibility are different rhetorical moves, and
proof-of-concept readers reasonably want both.

\subsection{The Anglo-Indian penal code family}
\label{subsec:anglo_indian_family}

The Singapore Penal Code 1871 is not a Singapore invention. It is the
Straits Settlements adoption of the 1860 Indian Penal Code drafted by
Macaulay's First Law Commission, lightly localised and subsequently
amended. The same parent text underpins the criminal codes of India
(IPC~1860, since 2023 reframed as the Bharatiya Nyaya Sanhita),
Pakistan (PPC~1860), Bangladesh (PC~1860), Brunei (PC, Cap.~22),
Malaysia (PC, Act~574), Sri Lanka (PC No.~2 of 1883), Myanmar
(PC~1861), and historically several East African codes. The shared
ancestry is not merely thematic. Section numbering, marginal-note
conventions, illustration formatting, exception clauses, and the
characteristic actus-reus / mens-rea / circumstance partitioning are
recognisably the same drafting style across the family. A reader who
has read Singapore's s415 will find the corresponding Indian and
Malaysian s415 cosmetically familiar, often differing only in
penalty quanta and amendment lineage.

This is what makes the family the natural test bed for grammar
portability. Yuho's primitives were chosen to fit the drafting shape
of penal statute --- not the cultural specifics of one jurisdiction.
The element partitioning is doctrinal, not local; the exception block
is a Catala-style prioritised override pattern that maps directly
onto the codified ``provided that'' clauses common to every member of
the family; the penalty combinators capture the cumulative-or-
alternative-or-conditional structure that the entire family inherits
from Macaulay's draft. To the extent the grammar fits Singapore
because the grammar fits Macaulay, the grammar should fit the others
for the same reason.

\subsection{What transfers, what adapts}
\label{subsec:portability_mechanism}

We separate the grammar substrate (which we expect to transfer
unchanged) from the per-jurisdiction configuration (which would need
swapping).

\paragraph{Transfers as-is.} The tree-sitter grammar, the AST, the
five-block statute structure (\texttt{elements}, \texttt{penalty},
\texttt{exceptions}, \texttt{illustrations}, \texttt{caselaw}), the
G10 reference resolver, the Z3 / Alloy verification scaffolding, the
seven transpilers (including the OASIS Akoma Ntoso target), and the entire editor surface (LSP, MCP, browser
extension, explorer site, Word add-in) are jurisdiction-neutral. They
operate on the AST, not on Singapore-specific tokens. Adopting a new
jurisdiction would not require parser changes for these subsystems.

\paragraph{Adapts via configuration.} A small number of surface
features are configurable per jurisdiction:

\begin{itemize}
\item \emph{Penalty vocabulary.} The set of admissible punishments
varies: caning is in the Singapore, Malaysian, and Bruneian codes
but not the Indian one; Pakistan retains \emph{tazir} categories
absent from Singapore; the Bangladeshi 2020 amendments enumerate
specific aggravators absent elsewhere. Yuho's punishment kinds are
a closed enum in the grammar; the per-jurisdiction adoption needs
the enum extended (or restricted) and the controlled-English
transpiler taught the new vocabulary.
\item \emph{Currency tokens.} Fines are encoded as a numeric quantum
plus a currency tag (\texttt{SGD} for Singapore). Indian, Pakistani,
and Bruneian codes use different currencies; the type system already
admits the parameter, but the L2 fidelity linter expects the
jurisdiction's reporting currency on initial scrape.
\item \emph{Statutory amendment lineage.} The \texttt{amends} clause
points at an act number (e.g.\ ``Act 32 of 2019''); each
jurisdiction has its own amendment-act numbering scheme. The grammar
accepts arbitrary strings here; only the L3 fidelity check that
reconciles \texttt{amends} against the canonical scrape needs a
jurisdiction-specific source.
\item \emph{Marginal-note style.} Older codes (1860--1900) place the
marginal note above the section; newer ones (post-1990) inline it.
The Phase~A scraper's heuristic for extracting marginal notes is
jurisdiction-specific; the AST and downstream tooling are not.
\end{itemize}

\paragraph{Requires substantive work.} Two surfaces are not yet
configuration-clean:

\begin{itemize}
\item \emph{Mens-rea taxonomy.} Singapore distinguishes intention,
knowledge, rashness, and negligence as distinct mental states; the
1860 IPC is identical, but the BNS~2023 (the recent Indian
recodification) introduces collapsed categories that do not map
one-to-one. A jurisdiction that uses a different taxonomy needs
either a grammar extension or a vocabulary remapping pass.
\item \emph{The L3 reviewer's checklist.} The 11-point fidelity
checklist (\S\ref{subsec:coverage}) bakes in conventions specific to
the Singapore SSO format (anchor IDs, marginal-note syntax,
amendment notation). A jurisdiction with a different official
publication format would need its own checklist; the
parameterisation surface exists but has not been exercised.
\end{itemize}

\subsection{Beyond the Penal Code family}
\label{subsec:non_penal}

The grammar is statute-shaped rather than penal-statute-shaped at the
top level: a module contains statutes, a statute contains elements
and penalties and exceptions, and the AST does not require the
penalties block to be present. The question of whether Yuho fits
non-penal statutes (tax, employment, intellectual property,
competition law) reduces to whether those statutes share the
\emph{element + exception + illustration} shape. The Singapore
Income Tax Act and the Employment Act both exhibit this shape on
inspection of a sample of sections; the Companies Act does too,
though with a heavier reliance on definitions blocks that Yuho
already supports. A more rigorous claim awaits empirical encoding.

The shift to civil-law jurisdictions is more speculative. The French
\emph{Code p\'enal} and the German \emph{Strafgesetzbuch} are
similarly element-structured, but their definitional density and
their reliance on cross-references to the civil-law general parts
mean an encoding would lean heavily on Yuho's reference graph and
less on its element partitioning. We do not claim the grammar fits
civil-law codes; we claim the AST is the right level at which to
test the question. Catala~\cite{catala2021}, focused on French tax
law, is the closer analogue here, and a Yuho-Catala interop layer
(scope composition over Yuho statute blocks) is the most natural
extension if the civil-law direction is pursued.

\subsection{An empirical path to validation}
\label{subsec:portability_path}

We commit to a concrete and falsifiable test of the portability
claim. Encoding a 50-section sample of the Indian Penal Code is the
minimum viable exercise: 50 sections is enough to surface grammar
gaps that the Singapore exercise did not, and small enough to
complete in a single encoder-quarter. The sampling strategy mirrors
the Singapore encoding order: start with structurally simple
property offences (theft, cheating, criminal misappropriation),
move through bodily-harm hierarchies (hurt, grievous hurt, culpable
homicide), and end with the procedural and abetment apparatus.
Concretely, the test passes if (a)~the grammar absorbs the 50-section
sample at L1 with zero new gaps catalogued, (b)~the L2 fidelity
linters produce diagnostic densities within an order of magnitude of
the Singapore baseline, and (c)~the existing G1--G14 gap catalogue
covers any drafting patterns the IPC sample exhibits but Singapore's
PC does not. Failure on (a) implies a fifteenth gap and a grammar
extension; failure on (b) implies the linters are over-fitted to
Singapore conventions; failure on (c) implies the gap catalogue is
not a stable artefact and would need a replication study before any
portability claim is publishable.

The corpus directory layout is already jurisdiction-aware: the slim
corpus, the canonical per-section JSON, and the reference graph are
keyed on a \texttt{jurisdiction} field that defaults to Singapore but
admits any string. Switching corpus base from Singapore to India
requires no code change to the editor surfaces, only a parallel
\texttt{library/penal\_code\_in/} directory and a per-jurisdiction
\texttt{config.toml} for the per-jurisdiction parameters above. The
build pipeline already supports running the corpus build over an
arbitrary library root, so a multi-jurisdiction artefact is an
engineering exercise on top of the empirical encoding work.

We frame this as future work in \S\ref{sec:conclusion} rather than as
a present-tense claim. What this section adds beyond
\S\ref{sec:limitations} is a positive structural argument: the
grammar substrate is jurisdiction-neutral, the configuration surface
is small, and the empirical test is bounded. Portability has not been
demonstrated, but the cost of demonstrating it is well-understood and
the path to falsification is clear.

\section{Limitations}
\label{sec:limitations}

\paragraph{Common-law doctrines are not modelled.} Yuho captures the
statute as drafted; it does not model judicial gloss. Doctrines like the
\emph{thin skull} rule or \emph{novus actus interveniens} live in case
law and shape how elements (especially causation) apply. Our \texttt{caselaw}
block records the existence of judicial interpretation but does not lift
the doctrine into the type system.

\paragraph{Single jurisdiction.} The encoding is Singapore-only. The
Anglo-Indian heritage of the Singapore Penal Code makes the grammar
plausibly portable to India, Pakistan, Brunei, and (with adjustments)
Malaysia, but we have not verified this empirically.

\paragraph{Shallow precedent integration.} Case law is referenced by name
+ citation + holding string; the holding is not parsed into structured
elements. A holding like ``mens rea may be inferred from the actus reus
when the act is inherently dangerous'' would, in a richer system, modify
the burden qualifier of a mens-rea element. Yuho does not yet do this.

\paragraph{No live SSO diffing.} The canonical scrape lives in
\texttt{library/penal\_code/\_raw/act.json} as a snapshot. Amendments
published by AGC are not picked up automatically. A scheduled re-scrape
+ diff pass is on the roadmap.

\paragraph{L3 reviewer is the encoder.} The 11-point checklist is
mechanical, but the same author wrote the encodings and stamped them.
External review by a Singapore-qualified lawyer is the natural next step
for any encoding intended for use beyond proof of concept.

\paragraph{Verification scope.} Z3 and Alloy hookups currently target
exception-priority well-formedness and bounded element-combination
enumeration. Properties that require quantification over an unbounded
number of historical events (e.g.\ ``no defendant has ever been convicted
under both s299 and s300 for the same act'') are out of scope.

\section{Future work}
\label{sec:future_work}

Three concrete strands, plus a longer roadmap of explicitly deferred
work.

\subsection{Concrete next strands}
\label{subsec:future_concrete}

The \textsf{G10} cross-section resolver landed during the most
recent iteration --- the LexScript-borrowed Tarjan-SCC analysis
(\S\ref{subsec:scc_eval}), the lam4-style \texttt{is\_infringed}
predicate, and the Catala-style \texttt{apply\_scope} composition
primitive are all shipped --- so the remaining work in this strand
is the elaboration of the evaluator-side semantics of
\texttt{apply\_scope} into a full element-graph executor that returns
bindings, not just a structural ``does this scope exist'' predicate.

\paragraph{Cross-jurisdiction generalisation.} 8 representative
Indian Penal Code 1860 sections (s300 / s302 / s375 / s376 / s378 /
s420 / s497 / s509) are encoded as a proof-of-concept demonstrating
grammar fit, with a fully-scraped IPC corpus (493 sections) and a
\texttt{yuho refs --compare-libraries} structural-overlap tool
shipped alongside. Full IPC encoding (493 sections) and a richer
SCC-overlap / amendment-divergence analysis are deferred to the
agent-dispatch pipeline. The structural argument and the
falsification criteria are laid out in \S\ref{sec:applicability};
the Akoma Ntoso transpiler is the on-ramp for any cross-jurisdiction
work that wants to interoperate with LegalDocML-consuming tooling.

\paragraph{Deeper precedent integration.} Lifting case-law holding
strings into structured constraints that modify element burden
qualifiers when the matching predicate fires. Right now the L3
stamp records caselaw anchors as opaque strings; the next step is
making them load-bearing.

\paragraph{Evaluator-side \texttt{apply\_scope}.} Elaborating the
cross-section composition primitive into a full element-graph
executor that returns bindings, not just a structural ``does this
scope exist'' predicate.

\subsection{Roadmap (Phase 2A / 2B / 2C)}
\label{subsec:future_roadmap}

The Yuho corpus and tooling extend along three orthogonal axes that
we treat as deferred but explicit, not unscoped.

\paragraph{Phase 2A --- diachronic depth.} The Penal Code 1871 has
150+ years of amendments. Phase 2A extends every encoded section
with every historical version since 1872; \texttt{yuho check
--as-of 1995-06-15 my\_facts.yh} becomes a first-class temporal
query. Acceptance: 524 sections $\times$ at-least-3 historical
versions encoded; per-section timeline rendering on
\texttt{/s/<N>.html}; novel paper contributions on encoding
lineage, an amendment-friction Levenshtein metric, and stable
kernels classifying sections by drift since 1872.

\paragraph{Phase 2B --- vision pivot.} Yuho moves beyond the Penal
Code into the wider Singapore criminal-justice corpus: Evidence
Act 1893, Criminal Procedure Code, Constitution 1963, and Contract
Law (common-law plus Misrepresentation Act and the Contracts
(Rights of Third Parties) Act). Civil remedies and equitable
doctrines need AST shapes Yuho does not yet have, so this is
gated on grammar extension first.

\paragraph{Phase 2C --- interactive learn-by-doing.} The static
explorer site grows into a guided learning surface for law
students: per-section JS fact-builders, in-browser evaluators,
worked-example walkthroughs from the canonical illustrations.
Co-authoring with a law-school educator who can run a classroom
trial unlocks measured learning-outcome evidence; the pedagogical
claim is then ``interactive structural pedagogy for common-law
penal codes''.

\medskip

Phases 2A--2C are deliberately out of scope for the present
artefact: 2A is a multi-month diachronic re-encode, 2B is a
multi-statute scope expansion, and 2C is a co-authored pedagogical
study. The current paper's empirical claims rest entirely on the
static, current-version, single-Act corpus; the roadmap is
provided for honest scoping rather than as load-bearing future
work.

\section{Conclusion}
\label{sec:conclusion}

We presented Yuho, a statute-shaped domain-specific language, and used
it to encode the entire Singapore Penal Code 1871: \statRawSections{}
sections, \statLOnePass{}/\statRawSections{} passing parse-and-build at
the L1+L2 coverage tiers, and \statLThreePass{} carrying a structured
human-fidelity stamp at the strictest L3 tier (\S\ref{subsec:coverage_eval}).
The exercise surfaced a catalogue of fourteen grammar gaps
(\S\ref{subsec:gaps}), distributed across ten genuine grammar
limitations that we resolved in the parser and four diagnostic-layer
issues that we resolved by adding fidelity linters. The result is a
complete, queryable, machine-checkable encoding of a production penal
code, plus the tooling --- a tree-sitter parser, an immutable Python
AST, eight transpilers (including Akoma Ntoso for cross-jurisdictional
interop and a Mermaid mindmap surface), a reference-graph resolver with Tarjan-SCC analysis and the
lam4-style \texttt{is\_infringed} / Catala-style \texttt{apply\_scope}
cross-section predicates, a Z3-and-Alloy verification hookup, a Language
Server, a Model Context Protocol server, and a VS Code extension ---
that turns the encoding into an executable artefact rather than a
paper exercise.

The longer-term thesis is that \emph{statute-as-source-of-truth} is a
viable posture for legal infrastructure. This is the
isomorphism~\cite{benchcapon1992isomorphism} position with all the
standard objections that come with it: clause-for-clause mirroring
preserves traceability and makes amendments propagate cleanly, but
it forecloses some logic-friendly normalisations and inherits the
open-textured-terms problem. We claim the isomorphism trade is the
right one for penal code; we do not claim the standard objections
have been dissolved. The encoded artefact is what
the lawyer, the tooling, and the verification stack all read;
downstream renderings (controlled English, \LaTeX{}, Mermaid, DOCX)
are generated, not authored. We do not claim that the \statRawSections{}
encoded sections constitute publication-ready legal advice; the
limitations on common-law-doctrine modelling, precedent integration,
and the author-administered L3 stamping (\S\ref{sec:limitations})
are explicit about what the artefact is and is not. We do claim that the gap
between a drafted penal section and an executable encoding is
narrower than the existing literature suggests, and that an
appropriately-shaped grammar dissolves a class of previously-reported
encoding obstacles, leaving a residue that is tractable at full
coverage for one statute family.

\paragraph{Artefact availability.} The Yuho compiler, encoded library,
Lean 4 mechanisation, and evaluation harness are available under an
open-source licence at \url{https://github.com/gongahkia/yuho}. A
versioned snapshot of this paper plus the v5.1.0 corpus is archived
at Zenodo, DOI
\href{https://doi.org/10.5281/zenodo.19935537}{10.5281/zenodo.19935537}.
The paper itself, including the build pipeline that injects coverage
statistics into this very manuscript at compile time, is in the same
repository under \texttt{paper/}.


\bibliographystyle{ACM-Reference-Format}
% Force numbered citations (acmart's default is author-year; the
% "nonacm" build wants numbered to match the screenshot format).
\setcitestyle{numbers,sort&compress}
% Force the Yuho artefact self-reference into the bibliography even
% though it is cited only contextually in the Artefact-availability
% paragraph; \nocite{yuho2026} guarantees the entry appears in the
% rendered references list.
\nocite{yuho2026}
\bibliography{references}

\appendix
\section{Scope of the formal treatment}
\label{app:semantics-full}

The operational rules in \S\ref{sec:semantics} cover the
core element-evaluation, exception-priority, and penalty
combinator judgments --- the fragment exercised by every
encoded section in the corpus and by the soundness lemmas
in \S\ref{sec:soundness}. Four refinements are deliberately
left at the prose level rather than fully formalised here:
mens-rea intent typing (the discrimination between
\emph{intentionally}, \emph{knowingly}, \emph{rashly}, and
\emph{negligently} as Yuho enum values, beyond the boolean
truth conditions of the present rules); the G7 temporal-
ordering constraint between actus reus and the reus-bearing
mens rea event; the G11 burden-of-proof qualifier on the
defendant side of an exception; and the cross-section
composition rules invoked when one section's elements
predicate over another section's offence. The Lean 4
mechanisation in \S\ref{sec:soundness}'s appendix discharge
covers G8/G14 sentinels and the conviction-layer oracle
constructively; the four refinements above remain as
inventory items for a v6 mechanisation pass and are
flagged in \S\ref{sec:limitations}.

\end{document}
